\documentclass{article}
\usepackage{filecontents}
\usepackage{pgfplotstable}
\usepackage{booktabs}

\begin{document}

% Table 1
\begin{table}[h]
\caption{Table 1: Current vs Voltage}
\centering
\pgfplotstabletypeset[
    col sep=comma,
    string type,
    columns/I/.style={column name={\textbf{I (mA)}}},
    columns/V/.style={column name={\textbf{V (V)}}},
    every head row/.style={before row=\toprule,after row=\midrule},
    every last row/.style={after row=\bottomrule}
]{data1.csv}
\end{table}

% Table 2
\begin{table}[h]
\caption{Table 2: Current vs Voltage}
\centering
\pgfplotstabletypeset[
    col sep=comma,
    string type,
    columns/I/.style={column name={\textbf{I (mA)}}},
    columns/V/.style={column name={\textbf{V (V)}}},
    every head row/.style={before row=\toprule,after row=\midrule},
    every last row/.style={after row=\bottomrule}
]{data2.csv}
\end{table}

\end{document}
